

\subsection{模型特点}
本文以统一物理约束和因果信息集贯通四个问题，仅按赛题改变合同调整权限与电价时序；实测真值仅在当期决策生成后进入执行修正和结算。预测、风险裕量、合同规划、滚动执行与现金结算构成闭环，所有方案均通过逐时段供需平衡、SOC递推、费用分项和互斥性检验。

问题三、四复用同一滚动调度框架，通过预报组合、边际收益和理想电价基准分别识别日内信息与电价未知的经济影响；同时报告平均费用、应急购电与尾部风险，避免以单一指标作出结论。六层提前期风险裕量保留了“预报越远、误差通常越大”的时序结构；对21:00的补充检验也表明，数值收益必须与新增信息的物理意义共同解释。

\subsection{局限与推广}
题目未给出并网容量上限以及售电相关机制，本文假设外购电无上限，且微网不可向主网反向送电。若工程场景存在并网容量约束，则需增加约束$g_t+u_t \le G_{\text{max}}$。模型采用经验分位数刻画历史残差尾部特征，未刻画连续多时段预测误差间的相关性，因此逐时段覆盖率不能等价为整日可靠概率。Q3、Q4的收益还可能同时来自信息更新、合同调整和重新优化，当前对照不能完全识别三者的边际贡献。仿真采用10分钟聚合时间粒度，忽略时隙内部测量延迟与功率爬坡过程；实际工程应用时，可缩短仿真步长或增加延迟状态变量予以修正。本文全部结论基于2025年历史数据得到，更换其他年份数据集时，需要重新校验模型参数与费用表现。

按预报提前时间分组设置风险裕量时，采用同一种风险裕量的全局方案，在样本上的总费用相比分组方案低8860.56元。本文仍然选用分组方案，原因在于该方案保留了“预报提前期越长，预测误差通常越大”的数据固有特征。问题三中$H=1$对应的样本费用略低于$H=6$，说明滚动时域长度需要在样本经济性与多时段规划能力之间做折中取舍。

模型可扩展到含风电、可控柔性负荷与多组储能的园区微网场景。新增分布式风电资源，只需在供需平衡方程内增加对应出力供给变量；引入需求响应资源，则增加负荷平移变量，并配套用户用电舒适度约束。若放开售电权限或存在负电价场景，应当采用带有显式充放电互斥约束的混合整数规划模型，规避充放电循环套利问题。微网分层控制的相关研究可参见\cite{bouzid2015microgrid}。

若能够获取多年实测数据，可将当前35日滑动分位替换为季节条件分位或场景树方法，并采用独立年份样本完成参数$W,\alpha,H$的优选；当前参数与$\lambda$均在同一年度样本上选择并评价，存在样本内调参偏差，正式应用时应采用时间滚动交叉验证或独立年份检验。如果预报供应商按次收取预报费用，可将$\kappa$纳入预报获取方案优化，把盈亏平衡分析拓展为0-1整数选择模型；当园区配备多块储能电池时，可以对每一块电池单独建立SOC状态方程，额外增加共享并网容量约束与电池寿命折损成本。上述所有拓展场景，均可沿用本文的购电计划生成、日内实际调度以及费用结算完整流程。
